\section{System Model}

Consider a cooperative downlink network with base-station set $\mathcal{B} = \{1,\dots,B\}$. Base station $b$ has $M$ transmit antennas and serves a user set $\mathcal{U}_b$. Let $\mathcal{K} = \cup_{b \in \mathcal{B}} \mathcal{U}_b$ denote the global user index set, and let $b(k)$ be the serving base station for user $k$.

Base station $b$ transmits
\begin{equation}
\vect{x}_b = \sum_{u \in \mathcal{U}_b} \vect{w}_{b,u} s_u,
\end{equation}
where $\vect{w}_{b,u} \in \C^M$ is the beamforming vector for user $u$ and $s_u$ is a unit-power information symbol. The received signal at user $k$ is
\begin{equation}
y_k = \sum_{b \in \mathcal{B}} \vect{h}_{k,b}^{\herm} \vect{x}_b + n_k,
\label{eq:received_signal}
\end{equation}
where $\vect{h}_{k,b} \in \C^M$ is the channel from base station $b$ to user $k$ and $n_k \sim \mathcal{CN}(0,\sigma_k^2)$.

With single-user decoding, the signal-to-interference-plus-noise ratio (SINR) for user $k$ is
\begin{equation}
\mathrm{SINR}_k(\mathbf{W}) =
\frac{\left|\vect{h}_{k,b(k)}^{\herm} \vect{w}_{b(k),k}\right|^2}
{\sum_{u \in \mathcal{U}_{b(k)} \setminus \{k\}} \left|\vect{h}_{k,b(k)}^{\herm} \vect{w}_{b(k),u}\right|^2
+ \sum_{c \neq b(k)} \sum_{u \in \mathcal{U}_c} \left|\vect{h}_{k,c}^{\herm} \vect{w}_{c,u}\right|^2
+ \sigma_k^2}.
\label{eq:sinr}
\end{equation}
The corresponding weighted rate is $\alpha_k R_k(\mathbf{W})$, where
\begin{equation}
R_k(\mathbf{W}) = \log_2 \bigl(1 + \mathrm{SINR}_k(\mathbf{W})\bigr).
\end{equation}

The globally coordinated benchmark solves
\begin{align}
\max_{\mathbf{W}} \quad & F(\mathbf{W}) = \sum_{k \in \mathcal{K}} \alpha_k R_k(\mathbf{W}) \label{eq:global_objective} \\
\text{s.t.} \quad &
\sum_{u \in \mathcal{U}_b} \norm{\vect{w}_{b,u}}^2 \le P_b, \quad \forall b \in \mathcal{B},
\label{eq:sum_power_constraint}
\end{align}
with optional per-antenna constraints represented by quadratic forms $\sum_{u \in \mathcal{U}_b} \vect{w}_{b,u}^{\herm} \mathbf{E}_{b,q} \vect{w}_{b,u} \le P_{b,q}$.

\subsection{Localization Model}

Let $d(b,c)$ denote the physical distance between base stations $b$ and $c$. For a radius $\rho \ge 0$, define the cooperation neighborhood
\begin{equation}
\mathcal{C}_b(\rho) = \left\{ c \in \mathcal{B} : d(b,c) \le \rho \right\}.
\end{equation}
Base station $b$ only exchanges optimization variables with stations in $\mathcal{C}_b(\rho)$. Accordingly, the localized objective retains interference couplings induced by nearby stations and truncates far-field interactions:
\begin{equation}
F^{(\rho)}(\mathbf{W}) = \sum_{k \in \mathcal{K}} \alpha_k R_k^{(\rho)}(\mathbf{W}),
\label{eq:localized_objective}
\end{equation}
where $R_k^{(\rho)}(\mathbf{W})$ is obtained from \eqref{eq:sinr} by retaining only those interference terms generated by base stations in $\mathcal{C}_{b(k)}(\rho)$. When needed for a conservative approximation, one may also include a residual term $\eta_k^{(\rho)}$ in the denominator to upper-bound omitted far-field interference. The resulting problem is a radius-dependent approximation to \eqref{eq:global_objective}--\eqref{eq:sum_power_constraint}.
